\documentclass[10pt,a4paper,twocolumn]{article}

\usepackage[utf8]{inputenc}
\usepackage[T1]{fontenc}
\usepackage{lmodern}
\usepackage[margin=1.8cm,columnsep=0.7cm]{geometry}
\usepackage{amsmath,amssymb,amsthm}
\usepackage[round,authoryear]{natbib}
\usepackage{xcolor}
\usepackage{hyperref}
\usepackage{enumitem}
\usepackage{booktabs}

\hypersetup{colorlinks=true,linkcolor=blue!60!black,
  citecolor=green!50!black,urlcolor=blue!70!black}

\theoremstyle{plain}
\newtheorem{theorem}{Theorem}
\newtheorem{proposition}[theorem]{Proposition}
\newtheorem{lemma}[theorem]{Lemma}
\newtheorem{corollary}[theorem]{Corollary}
\newtheorem{conjecture}[theorem]{Conjecture}
\theoremstyle{definition}
\newtheorem{definition}{Definition}
\newtheorem{remark}{Remark}

\newcommand{\dd}{\mathrm{d}}
\newcommand{\ee}{\mathrm{e}}
\newcommand{\eps}{\varepsilon}
\newcommand{\Cstr}[1]{\mathcal{C}_{#1}}
\newcommand{\Tree}{\tau}
\newcommand{\USM}{\mathrm{UST}}
\newcommand{\LERW}{\mathrm{LERW}}
\newcommand{\Tr}{\operatorname{tr}}
\newcommand{\Ex}{\mathbb{E}}

\usepackage{mdframed}
\newmdenv[backgroundcolor=blue!5,linecolor=blue!40,
          linewidth=1pt,innertopmargin=6pt,innerbottommargin=6pt,
          innerleftmargin=8pt,innerrightmargin=8pt]{slotbox}

\title{Counting Past the Boundary\\[2pt]
  \large Enumerative CFAC for LERW in $d{=}3$, TAP Complexity, and KPZ Upper Tails}

\author{S. Amarteifio\\
  \small\textit{Percolation Labs}}
\date{April 2026}

\begin{document}
\maketitle

\begin{abstract}
The CFAC stratification theorem~\citep{AmarteifioCFAC2026} has a narrow
domain: polynomial Dyson--Schwinger equations with algebraic or
$D$-finite generating functions and a scalar rank-$k$ bridge.  The
companion paper~\citep{AmarteifioMannaC3} marks three cases beyond
that boundary --- loop-erased random walk in three dimensions, TAP
complexity in spin glasses, and the KPZ upper-tail rate function.
Here we revisit those cases under the broader CFAC umbrella, where
the theorem is one deliverable and \emph{weighted enumeration of
combinatorial objects} is another.  For each case we extract a clean
generating-function statement: (i) the LERW two-point function on any
finite graph is a ratio of Kirchhoff polynomials, so the 3D exponent
$\nu$ is a finite-size scaling question about a concrete rational
function; (ii) the Kac--Rice expected complexity for $p$-spin models
is (up to a Legendre transform) a ribbon-graph enumeration, evaluated
via Harer--Zagier; (iii) the replicated lattice directed polymer is a
chemical reaction network of $n$ distinguishable walkers with
pairwise $\delta$-attraction, with Schur-process structure.  None of
these produces the open exponent in closed form, but each gives the
object under study a concrete combinatorial definition that the
CFAC pipeline can operate on at finite size.  The purpose is to
widen the scope statement honestly: the theorem does not reach these
problems, but the enumeration layer does.
\end{abstract}

\tableofcontents

% ==================================================================
\section{Introduction: two tiers of CFAC}
\label{sec:intro}

The CFAC programme, as stated in~\citep{AmarteifioCFAC2026}, has two
distinct deliverables that are usually discussed together.

\begin{description}
\item[Tier A (stratification theorem).] For a polynomial
  Dyson--Schwinger equation $G = z\phi(G)$ with an algebraic or
  $D$-finite generating function, the Puiseux order $k$ of the
  dominant branch point gives a density-exponent skeleton
  $\tau_0 = 1 + 1/k$, with a one-loop dressing $\gamma_k(d)$ from a
  scalar rank-$k$ bridge function.
\item[Tier B (enumerative combinatorics).] The underlying CFAC
  pipeline treats Feynman diagrams as weighted combinatorial objects
  --- Symanzik polynomials, matroidal deletion--contraction,
  BPHZ multiplicities, sector decomposition --- and computes exact
  finite-size generating functions independently of whether a
  closed-form asymptotic exists.
\end{description}

\subsection*{Background: unpacking the Tier-A statement}
\label{sec:intro-bg}

The Tier-A sentence compresses several technical notions that are
worth spelling out for readers coming from physics or probability
rather than analytic combinatorics.

\paragraph{Generating function (GF).}
A formal power series $F(z) = \sum_{n \geq 0} a_n z^n$ whose
coefficients $a_n$ count the combinatorial configurations of size
$n$ in the problem of interest --- for us, Feynman diagrams with
$n$ loops, walks of length $n$, spanning trees of given weight,
etc.  The variable $z$ has no physical meaning; it is a
book-keeping marker.  For a physicist, $F(z)$ plays the role of
a partition function in a fictitious ensemble with fugacity $z$
for ``size''.  The asymptotic behaviour of $a_n$ as $n \to \infty$
--- which is where physics lives --- is encoded in the nearest
singularity of $F(z)$ on the circle of convergence: location
$z_c$ sets the exponential growth $a_n \sim z_c^{-n}$, and the
\emph{type} of singularity sets the subleading power.

\paragraph{Algebraic GF.}
$F(z)$ is \emph{algebraic} if it satisfies a polynomial equation
$P(z, F(z)) = 0$ for some $P \in \mathbb{Q}[z, y]$.  The simplest
example is the Catalan-number GF
$F = 1 + z F^2$, solution $F(z) = (1 - \sqrt{1 - 4z}) / (2z)$,
with a square-root branch point at $z_c = 1/4$.  Algebraic means
``closed-form in polynomials and $n$-th roots''.  Many tree-like
combinatorial classes have algebraic GFs; all of them admit exact
asymptotic extraction by singularity analysis.

\paragraph{$D$-finite GF.}
$F(z)$ is \emph{$D$-finite} (equivalently, \emph{holonomic}) if it
satisfies a linear ordinary differential equation with polynomial
coefficients:
$p_r(z) F^{(r)}(z) + \cdots + p_0(z) F(z) = 0$
with $p_i \in \mathbb{Q}[z]$.  Every algebraic GF is $D$-finite,
but the converse fails: $\exp(z)$, Bessel functions,
hypergeometric series, and solutions of most special-function
ODEs are $D$-finite but not algebraic.  For a physicist, $D$-finite
is the class where the function is ``solvable by a finite
ODE with polynomial coefficients'', and its coefficients satisfy
a finite linear recurrence with polynomial coefficients.
\emph{Not} $D$-finite means: no such finite ODE exists; the
function lies outside the reach of classical special-function
methods.  Lattice Green's functions on $\mathbb{Z}^d$ for
$d \geq 2$ are the canonical non-$D$-finite examples in physics.

\paragraph{Dyson--Schwinger equation (DSE).}
In quantum field theory a DSE is a self-consistent integral
equation relating a full Green's function to itself, obtained by
resumming an infinite tower of Feynman diagrams.  In our
combinatorial setting, ``DSE'' specialises to a functional
equation of the form
\[
  G(z) \;=\; z \, \phi\bigl(G(z)\bigr),
\]
where $\phi$ is a polynomial (or more general analytic function)
of $G$.  This is the ``polynomial DSE'' of Tier A.  It encodes
rooted tree-like combinatorial structures: the root contributes
one $z$, each child subtree contributes a factor of $G$, and the
vertex structure lives in $\phi$.  The Catalan recursion
$G = z(1 + G^2)$ (rooted binary trees) and more generally any
vertex polynomial give specific instances.  Self-similar
random-cluster models, directed-percolation mean-field
recursions, and Doi--Peliti tree-level vertex resummations all
produce DSEs of this shape.

\paragraph{Puiseux order $k$ at the dominant singularity.}
Near its nearest singularity $z_c$ an algebraic GF admits a
\emph{Puiseux expansion} --- a power series in fractional powers
of $(z_c - z)$:
\[
  G(z) \;=\; G(z_c) \;-\; A\,(z_c - z)^{1/k} \;+\; \cdots,
  \qquad k \in \mathbb{Z}_{\geq 2}.
\]
The integer $k$ is the \emph{Puiseux order}.  For $k = 2$ the
leading behaviour is a square-root cusp, for $k = 3$ a cube-root
cusp, and so on.  By singularity analysis the coefficients behave
as $a_n \sim C z_c^{-n} n^{-1 - 1/k}$; the exponent
$\tau_0 = 1 + 1/k$ appearing in Tier A is this.  Different
universality classes correspond to different $k$: mean-field
directed percolation has $k = 2$, tricritical branching has
$k = 3$, and so on.  The Puiseux order $k$ is the
\emph{CFAC stratification index}: a discrete invariant separating
the space of polynomial DSEs into strata, within each of which
the leading exponent $\tau_0$ is universal.

\paragraph{Scalar rank-$k$ bridge.}
Beyond the leading exponent $\tau_0 = 1 + 1/k$, the CFAC
programme refines the universal answer with a one-loop correction
$\gamma_k(d)$ encoding the spatial dimension $d$.  The object
that carries this correction is a \emph{bridge function}
$B_k(r)$: a universal single-variable function of a dimensionless
ratio $r$ (mass or diffusion ratio, etc.) determined by the rank
$k$ and independent of the specific DSE.  ``Scalar'' means $B_k$
is number-valued (the bridge is a scalar function of one
argument); ``rank-$k$'' refers to the Puiseux order $k$ that
classified the DSE.  The bridge converts the counting exponent
$\tau_0$ into a physical density exponent, with explicit
dimensional dependence $\gamma_k(d)$ inserted at the one-loop
order.  The failure modes in the three boundary problems of the
companion paper~\citep{AmarteifioMannaC3} are precisely failures
of this scalar bridge to close: the problem demands a bridge of
higher tensor rank (TAP), a disorder-averaged bridge (KPZ upper
tail), or there is simply no DSE of the polynomial form at all
(LERW in $d = 3$ --- the GF is non-$D$-finite).

\paragraph{Stratification.}
Putting these together: the CFAC theorem \emph{stratifies} the
space of polynomial DSEs by the Puiseux order $k$ of the dominant
singularity, and attaches to each stratum a canonical pair
$(\tau_0, \gamma_k)$ --- the skeleton exponent and the one-loop
dressing from the scalar rank-$k$ bridge.  The word
``stratification'' is used in the classical sense (decomposition
of a space into lower-dimensional pieces indexed by a discrete
invariant), not in any sheaf-theoretic sense.

\medskip

The companion paper~\citep{AmarteifioMannaC3} flags three problems
where Tier A does not reach:

\begin{itemize}[leftmargin=*]
\item LERW in $d = 3$: the thermodynamic-limit generating function on
  $\mathbb{Z}^3$ is a lattice Green's function, neither algebraic nor
  $D$-finite.
\item TAP complexity: the organising quantity $\Ex|\det(H - fI)|$
  needs a matrix-valued bridge; the current bridge is scalar.
\item KPZ upper tail: the coefficient of the $|h|^{3/2}$ rate
  function requires a disorder-averaged bridge; current bridges are
  deterministic.
\end{itemize}

Each obstruction is a Tier-A obstruction.  In this note we ask the
Tier-B question for the same three problems: \emph{what weighted
enumeration does the generating-function language hand us}?  The
answers are known to specialists in each subfield; the contribution
here is to show that they align under a single CFAC lens, giving
concrete objects the existing machinery can compute at finite size.

Throughout, ``CFAC enumeration'' means the diagrammatic /
combinatorial content of the pipeline --- Symanzik polynomials,
matrix-tree determinants, weighted path and partition sums --- with
no appeal to the stratification theorem.

% ==================================================================
\section{The enumerative layer}
\label{sec:enumerative}

Three constructions recur below.  Fixing them here lets the three case
sections stay short.

\paragraph{Kirchhoff / first Symanzik polynomial.}  For a finite
graph $G = (V, E)$ with edge weights $w : E \to \mathbb{R}_{\geq 0}$,
the \emph{weighted spanning-tree polynomial} is
\begin{equation}
  \Tree(G; w) \;=\; \sum_{T \in \mathcal{T}(G)} \prod_{e \in T} w_e,
  \label{eq:kirchhoff}
\end{equation}
where $\mathcal{T}(G)$ is the set of spanning trees.  The
\emph{matrix-tree theorem} gives a closed form:
$\Tree(G; w) = \det L^{(v)}(G; w)$ for any reduced weighted Laplacian
$L^{(v)}$.  In parametric Feynman integration,
Eq.~\eqref{eq:kirchhoff} is the first Symanzik polynomial; in CFAC
it is the object produced by \texttt{rdft/ac/dse.py} after
matroidal deletion--contraction.

\paragraph{Map / ribbon-graph enumeration.}  For an
$N \times N$ symmetric Gaussian random matrix $M$ (GOE),
\begin{equation}
  \Ex[\Tr M^{2k}] \;=\; \sum_{g \geq 0} N^{k+1-2g} \, \varepsilon_g(k),
\end{equation}
where $\varepsilon_g(k)$ is the number of rooted one-face maps with $k$
edges on a (possibly non-orientable) surface of genus $g$.  The
orientable-surface case is the Harer--Zagier
formula~\citep{HarerZagier1986}; the non-orientable extension is due
to~\citep{GouldenJackson1997}.  In CFAC these are the Feynman
diagrams of the matrix-model action $\tfrac12 \Tr M^2 - \tfrac{g}{N}
\Tr M^k$, computed as ordinary (weighted) graph counts.

\paragraph{Weighted path enumeration with $\delta$-interactions.}
Fix a transition kernel $p(x, y)$ on a finite state space.  The
partition function of $n$ distinguishable walkers of length $N$
with a pairwise attractive contact interaction of strength $\beta^2$ is
\begin{equation}
  Z_N^{(n)}(\beta) \;=\; \sum_{\omega^1, \ldots, \omega^n} \prod_{t=1}^N
    p(\omega^{\cdot}_{t-1}, \omega^{\cdot}_t) \,
    e^{(\beta^2/2)\, k_t(\omega)},
  \label{eq:nwalker}
\end{equation}
where $k_t(\omega) = \sum_x (\#\{i : \omega^i_t = x\})^2$ is the
coincidence count at time $t$.  The reaction-network reading is a
CRN with $n$ species and a symmetric binary vertex
$i + j \to i + j$ carrying rate $\beta^2$.  Doi--Peliti quantisation
of this CRN gives the replica moment of a directed polymer
(§\ref{sec:kpz}).

Each case below reduces to one of these three constructions.

% ==================================================================
\section{LERW in $d = 3$ as Kirchhoff enumeration}
\label{sec:lerw}

\paragraph{Setup.} Let $G = (V, E)$ be a finite connected weighted
graph and fix two vertices $a, b \in V$.  The loop-erased random walk
$\LERW_{a\to b}$ is the law of the chronologically loop-erased
simple random walk from $a$ conditioned to hit $b$.  By Wilson's
algorithm~\citep{Wilson1996}, $\LERW_{a \to b}$ coincides with the law
of the unique $a$-to-$b$ path in the uniform spanning tree of $G$.

\begin{proposition}[LERW as spanning-tree ratio]
\label{prop:lerw}
For any simple path $\gamma$ from $a$ to $b$ in $G$,
\begin{equation}
  \mathbb{P}\bigl[\LERW_{a\to b} = \gamma\bigr]
  \;=\; \frac{w(\gamma) \, \Tree(G / \gamma)}{\Tree(G)},
  \label{eq:lerw-kirchhoff}
\end{equation}
where $w(\gamma) = \prod_{e \in \gamma} w_e$ and $G / \gamma$ is the
contraction of $G$ along $\gamma$.
\end{proposition}

\begin{proof}[Proof sketch]
The event $\{\LERW_{a\to b} = \gamma\}$ is the event that the unique
$a$-$b$ path in the UST equals $\gamma$, which is the event
$\{\gamma \subset \USM\}$.  The conditional law of the UST given
$\gamma \subset \USM$ is uniform over spanning trees of $G/\gamma$.
The unconditional UST weight of a fixed spanning tree $T$ is
$w(T)/\Tree(G)$, so
$\mathbb{P}[\gamma \subset \USM] = w(\gamma) \Tree(G/\gamma)/\Tree(G)$.
\end{proof}

\paragraph{Length generating function.}
Summing Eq.~\eqref{eq:lerw-kirchhoff} against a length marker gives a
rational function in $z$:
\begin{equation}
  F_{ab}(z; G) \;\equiv\;
  \sum_{\gamma : a \to b \,\text{simple}} z^{|\gamma|} \, w(\gamma) \,
  \frac{\Tree(G / \gamma)}{\Tree(G)}.
  \label{eq:Fab}
\end{equation}

\begin{corollary}[Rational representation on finite graphs]
\label{cor:rational}
For any finite $G$, $F_{ab}(z; G)$ is a rational function in $z$ with
denominator $\Tree(G)$.  Both numerator and denominator are
evaluations of the weighted Laplacian determinant.
\end{corollary}

The corollary is just bookkeeping on
Eq.~\eqref{eq:lerw-kirchhoff}, but the content is that the LERW
two-point function is a quantity CFAC already computes: the ratio of
two first Symanzik polynomials.

\paragraph{Scaling limit as finite-size scaling of
$\boldsymbol{F_{0b}}$.}
Take $G = \mathbb{Z}^3_N$ (the $N^3$-torus).  Let
$\langle |\gamma| \rangle_{0b}$ denote the LERW expected length from
$0$ to $b$ under Eq.~\eqref{eq:lerw-kirchhoff}.  The 3D LERW scaling
exponent~\citep{Kozma2007} can be read off from
\begin{equation}
  \langle |\gamma| \rangle_{0b} \;\asymp\; |b|^{1/\nu_3},
  \qquad \nu_3 \approx 1.624,
\end{equation}
which by Corollary~\ref{cor:rational} is a finite-size scaling
statement about $z \partial_z \log F_{0b}(z; \mathbb{Z}^3_N)\big|_{z=1}$
as $N, |b| \to \infty$ in the appropriate order.  The obstruction the
stratification theorem meets --- that the thermodynamic limit is a
non-$D$-finite lattice Green's function --- does not obstruct this
finite-size object, which is a ratio of explicit determinants.

\paragraph{Tropical / monomial-dominance limit.}
In the limit where one spanning tree dominates (e.g.\ when edge
weights are $e^{-\beta E_e}$ and $\beta \to \infty$),
Eq.~\eqref{eq:lerw-kirchhoff} reduces to a deterministic path: the
unique minimum-weight spanning tree.  The LERW distribution
collapses onto the tree-path in this limit.  This is the tropical
Kirchhoff limit mentioned in~\citep[\S7.1]{AmarteifioMannaC3}; it is
a well-defined CFAC enumeration in its own right.

\begin{remark}
Prop.~\ref{prop:lerw} is classical (it is Wilson's
algorithm in reverse).  The CFAC reading is that
Eq.~\eqref{eq:lerw-kirchhoff} places LERW squarely inside the
first-Symanzik-polynomial enumeration tier: the same machinery that
produces two-loop hat diagrams~\citep{Amarteifio2026LDWC} produces
the LERW two-point function, with no changes beyond the graph on
which it is evaluated.
\end{remark}

\paragraph{Why $\boldsymbol{d = 3}$ is the hard case: dimensional
partitioning of the LERW exponent.}
The exponent $\nu_d$ is tractable in every dimension \emph{except}
the physically relevant one, and the reason is structural:
\begin{itemize}[leftmargin=*]
\item \textbf{$d = 2$: conformal symmetry.}  The LERW scaling limit
  is SLE$_2$~\citep{LawlerSchrammWerner2004}, an instance of the
  conformally-invariant family SLE$_\kappa$.  Conformal invariance
  in two dimensions is generated by an infinite-dimensional Lie
  algebra (Virasoro), enough symmetry to collapse every observable
  to a formula in the single parameter $\kappa$.  For LERW this
  gives $d_f = 1 + \kappa / 8 = 5/4$ in closed form, and with it
  $\nu_2 = 4/5$.

\item \textbf{$d \geq 4$: mean-field / dimensional decoupling.}
  Above the upper critical dimension $d_c = 4$, two independent
  simple random walks typically do not intersect, so loop-erasure
  acts trivially in the bulk.  LERW behaves as SRW with at most
  logarithmic corrections, and $d_f = 2$ is forced by the Gaussian
  scaling of SRW~\citep{Lawler1986}.  At $d = 4$ this is exact up
  to $\log^{1/3}$ corrections; for $d \geq 5$ it is exact.

\item \textbf{$d = 3$: no symmetry, no small parameter.}  In three
  dimensions the conformal group reduces to a finite-dimensional
  M\"obius group (insufficient to fix any exponent), and the
  interaction is not small (two 3D random walks intersect with
  positive probability).  Kozma's theorem~\citep{Kozma2007} proves
  the scaling limit exists and is supported on a set of some
  Hausdorff dimension $d_f^{(3)}$, numerically
  $\approx 1.6236$~\citep{Wilson2010}, but the proof is
  non-constructive --- no technique is known that produces
  $d_f^{(3)}$ in any form, closed or otherwise.  The same
  obstruction is responsible for the unresolved status of 3D
  Ising, 3D percolation, and 3D Heisenberg critical exponents.
\end{itemize}

The obstruction is therefore \emph{dimension-specific} and
\emph{tool-specific}: 2D has conformal algebra, $d \geq 4$ has
dimensional decoupling, 3D has neither.  The finite-graph
Kirchhoff ratio of Prop.~\ref{prop:lerw} is exact in every
dimension; the scaling-limit extraction loses closed-form
tractability only at $d = 3$.  Section~\ref{sec:fss3d} returns
to this from the analytic-combinatorics angle: given that the
finite-size generating function is known exactly, what does
asymptotic extraction from it give us, even without a closed
form for the limit?

% ==================================================================
\section{TAP complexity as ribbon-graph enumeration}
\label{sec:tap}

\paragraph{Setup.}  For the spherical pure $p$-spin Hamiltonian
$H_N$ on $S^{N-1}(\sqrt{N})$, the Kac--Rice formula for the expected
number of critical values in $[NE_1, NE_2]$ is
\begin{equation}
  \Ex N_{\text{crit}}([E_1, E_2])
  \;=\; c_{N,p} \int_{E_1}^{E_2}
    \Ex \bigl|\det(M - f(E) I)\bigr| \, \rho_E(E) \, \dd E,
  \label{eq:kacrice}
\end{equation}
where $M$ is a GOE$(N{-}1)$ matrix, $f(E)$ is a linear shift
fixed by the Lagrange multiplier for the spherical constraint, and
$c_{N,p}, \rho_E$ are explicit Gaussian factors~\citep{ABAC2013,Subag2017}.

\paragraph{Characteristic-polynomial expectation as a generating
function.}
For GOE$(N)$,
\begin{equation}
  \log \Ex \bigl|\det(M - fI)\bigr|
  \;=\; \sum_{k \geq 1} \frac{(-1)^{k+1}}{k f^k}\,
        \Ex \Tr \bigl[(M/\sqrt{N})^k\bigr]
        + O(1/f),
  \label{eq:logdet-exp}
\end{equation}
convergent for $|f|$ above the semicircle edge.  The coefficients
$\Ex\Tr[(M/\sqrt{N})^k]$ admit exact finite-$N$ expansions into
rooted map counts~\citep{HarerZagier1986,GouldenJackson1997}:
\begin{equation}
  \Ex \bigl[\Tr (M/\sqrt{N})^{2k}\bigr]
  \;=\; \sum_{g \geq 0} N^{-2g}\, \varepsilon^{\rm GOE}_g(k),
  \label{eq:HZ-GOE}
\end{equation}
where $\varepsilon^{\rm GOE}_g(k)$ counts rooted one-face maps with
$k$ edges on a (possibly non-orientable) surface of genus $g$.

\begin{proposition}[TAP complexity as map enumeration]
\label{prop:tap}
The Kac--Rice complexity integrand in Eq.~\eqref{eq:kacrice} is the
exponential of a generating function over rooted one-face maps:
\begin{equation}
  \log \Ex \bigl|\det(M - fI)\bigr|
  \;=\; N \, \Phi_0(f) + \Phi_1(f) + N^{-2}\Phi_2(f) + \cdots,
\end{equation}
with each $\Phi_g$ determined by the map-counting coefficients
$\{\varepsilon^{\rm GOE}_g(k)\}_{k \geq 1}$ via
Eq.~\eqref{eq:logdet-exp} and the moment expansion
Eq.~\eqref{eq:HZ-GOE}.  At large $N$, $\Phi_0(f) = \int \log|\lambda -
f|\,\rho_{\rm sc}(\lambda)\,\dd\lambda$, recovering the semicircle
complexity density of~\citep{ABAC2013}.
\end{proposition}

\begin{proof}
Cluster expansion of $\log \Ex\,\exp[\sum_k \lambda_k \Tr M^k]$ against
the genus expansion of~\citep{HarerZagier1986, GouldenJackson1997};
specialise $\lambda_k$ to the coefficients of $\log|\det(M - fI)|$
from Eq.~\eqref{eq:logdet-exp}.  The large-$N$ limit of
$\Phi_0$ is the Wigner semicircle transform because the
planar ($g = 0$) map count is the Catalan numbers, whose generating
function is $(1 - \sqrt{1 - 4t})/(2t)$.
\end{proof}

\paragraph{CFAC hookup.}  $\Phi_0$ is the standard TAP complexity
density.  Prop.~\ref{prop:tap} places it in the CFAC enumeration
tier: the finite-$N$ and subleading corrections are exact sums over
rooted one-face maps, which are diagrams the pipeline handles
natively (they are the Feynman diagrams of a matrix $\phi^3 / \phi^4$
model).  The obstruction flagged
in~\citep[\S7.2]{AmarteifioMannaC3} --- that a matrix-valued bridge
is needed to close the scaling of $\Sigma(E)$ via the stratification
theorem --- does not apply at this tier: one is counting maps, not
extracting a Puiseux exponent.

\begin{remark}
The Legendre transform
$\Sigma(E) = \sup_f [\,-R(E, f) + \Phi_0(f)\,]$, with $R(E, f)$
the Gaussian rate~\citep{ABAC2013}, is the large-deviation form of
TAP complexity.  In the map-counting language, this is the saddle
point of a generating function over rooted ribbon graphs with an
external parameter $f$ tuned to the complexity curve.  The complexity
threshold $E_\infty$ (below which critical points cease to be
exponentially numerous) is the value of $E$ at which the saddle hits
the edge of the semicircle, i.e.\ the radius of convergence of the
planar map generating function.
\end{remark}

% ==================================================================
\section{KPZ upper tail as replica CRN enumeration}
\label{sec:kpz}

\paragraph{Setup.}  The $(1{+}1)$-dimensional lattice directed
polymer with i.i.d.\ Gaussian weights $\eta(t, x)$ of variance $1$ has
partition function
\begin{equation}
  Z_N(\beta) \;=\; \sum_{\omega : (0,0) \to (N, \cdot)}
    \prod_{t=0}^{N} \exp\!\bigl[\beta \eta(t, \omega_t)\bigr],
\end{equation}
with sum over nearest-neighbour lattice paths.  The KPZ universality
hypothesis predicts that for large $N$, $\log Z_N - c_1 N \sim c_2
N^{1/3} \chi$ with $\chi$ Tracy--Widom-distributed, and an upper-tail
rate function
$\mathbb{P}[\log Z_N - c_1 N > y N^{1/3}] \asymp e^{-c_3 y^{3/2}}$
at large positive $y$~\citep{CorwinGhosal2020, Krajenbrink}.

\paragraph{Replicated moments as weighted enumeration.}
For integer $n \geq 1$, disorder-averaging gives
\begin{equation}
  \Ex[Z_N(\beta)^n]
  \;=\; \sum_{\omega^1, \ldots, \omega^n}
        \prod_{t=0}^N
        \exp\!\Bigl[\tfrac{\beta^2}{2}\,
           k_t(\omega^{1:n})\Bigr],
  \label{eq:replica-moment}
\end{equation}
with the coincidence count
$k_t(\omega^{1:n}) = \sum_x (\#\{i : \omega^i_t = x\})^2$.
Eq.~\eqref{eq:replica-moment} is an instance of the weighted
$n$-walker enumeration Eq.~\eqref{eq:nwalker}.

\begin{proposition}[KPZ replica moments as CRN partition function]
\label{prop:kpz-replica}
For each integer $n \geq 1$, $\Ex[Z_N^n]$ is the partition function
of a chemical reaction network with $n$ distinguishable random-walker
species and a symmetric binary contact vertex
$i + j \to i + j$ of rate $\beta^2$, on $N$ time steps.  The
Doi--Peliti representation is an MSR action with kinetic term for
each walker and a single $\delta$-localised four-leg vertex
identifying species pairwise.
\end{proposition}

The proof is direct inspection of Eq.~\eqref{eq:replica-moment}.  In
the continuum limit the CRN becomes the attractive $n$-body
Lieb--Liniger model, and this is the classical replica route to
Tracy--Widom~\citep{Krajenbrink}.

\paragraph{Schur-process form (discrete geometries).}
For several exactly-solvable lattice models (log-gamma polymer,
O'Connell--Yor, geometric last passage), RSK correspondence converts
the replica sum into a partition sum~\citep{BorodinCorwin2014}:
\begin{equation}
  \Ex[Z_N^n] \;=\; \sum_{\lambda \vdash n} s_\lambda(x_1, \ldots, x_N)\,
      s_\lambda(y_1, \ldots, y_N) \cdot (\text{spectral weight}),
  \label{eq:schur-form}
\end{equation}
with $s_\lambda$ the Schur polynomial.  Eq.~\eqref{eq:schur-form} is
an enumeration over partitions weighted by Schur pairs, and the
generating function $\sum_n \lambda^n / n!\, \Ex[Z_N^n]$ becomes a
Schur measure on partitions --- a determinantal point process whose
right edge, under an appropriate scaling, is the Airy$_2$ kernel
underlying Tracy--Widom.

\paragraph{Upper-tail exponent as a Schur-measure edge statement.}
The $|y|^{3/2}$ upper-tail shape arises as the large-deviation form
of the right edge of the Schur / Plancherel measure on partitions,
an instance of the edge large-deviation theory of determinantal
processes.  This is where the \emph{coefficient} problem sits:
translating between lattice weight distributions and Schur specialisations.

\begin{proposition}[Exponent vs.\ coefficient]
\label{prop:kpz-exponent-coeff}
At the exponent level, the $|y|^{3/2}$ upper-tail shape follows from
any discretisation admitting the Schur-process resummation
Eq.~\eqref{eq:schur-form} together with Airy-kernel edge asymptotics
of the corresponding Schur measure.  The model-specific coefficient
$c_3$ depends on the spectral weight and is not determined by the
enumeration alone.
\end{proposition}

\begin{proof}[Proof sketch]
Given Eq.~\eqref{eq:schur-form}, the generating function
$\mathcal{G}(\lambda) = \sum_n \lambda^n \Ex[Z_N^n] / n!$ is a
Fredholm determinant for a Schur kernel~\citep{BorodinCorwin2014}.
The upper-tail rate of $\log Z_N$ corresponds to the rate of decay
of $\mathbb{P}[\text{top row of }\lambda > L]$ under the Schur
measure.  The classical edge large-deviation result for such measures
gives a rate function cubic in the scaled variable, yielding the
$|y|^{3/2}$ shape under the appropriate exponentiation~\citep{Krajenbrink}.
The coefficient depends on the derivative of the spectral density at
the edge, which is a model input.
\end{proof}

\paragraph{CFAC hookup.}  The CRN of Prop.~\ref{prop:kpz-replica}
sits inside the CFAC pipeline without modification: it is a Doi--Peliti
action with $n$ species and one vertex.  The enumeration of
configurations at fixed $n$ is a standard walker counting.  What is
\emph{not} inside the current pipeline is the replica continuation
$n \to 0$ and the disorder-averaged bridge that would fix $c_3$ from
first principles.

\paragraph{Tier-A closed form for the 2-body binding on
$\boldsymbol{\mathbb{Z}}$.}
At $n = 2$ the CRN reduces to a tight-binding chain with an
attractive on-site impurity, which admits a closed-form bound-state
eigenvalue.  This is the one place where the enumeration layer
delivers an exact Tier-A statement for the boundary cases.

\begin{theorem}[Closed-form 2-body replica rate]
\label{thm:d2-closed}
For the lattice directed polymer on $\mathbb{Z}$ with Gaussian
disorder of variance $1$ (coincidence weight
$\exp(\tfrac{\beta^2}{2} k_t)$, $k_t = \sum_x (\#\{i: \omega^i_t = x\})^2$),
the 2-body replica rate satisfies
\begin{align}
  D(2, \beta; \mathbb{Z})
  &\equiv \lambda(2, \beta; \mathbb{Z}) - 2 \lambda(1, \beta; \mathbb{Z})
     \nonumber \\
  &= 2\beta^2 - \log\bigl(2 \ee^{\beta^2} - 1\bigr).
  \label{eq:d2-closed}
\end{align}
The bound-state wavefunction in the relative coordinate $r = x_1 -
x_2$ is $\psi(r) = z^{|r|/2}$ with $z = 1/(2\ee^{\beta^2} - 1)$.
\end{theorem}

\begin{proof}
Separate the translation-invariant centre-of-mass coordinate
$s = (x_1 + x_2)/2$ from the relative coordinate $r$.  At
$k_s = 0$ the relative transfer matrix acts as
$(T_{\text{rel}}\psi)(r) = \psi(r-2) + 2\psi(r) + \psi(r+2)$ for
$r \neq 0$, with a multiplicative $\ee^{\beta^2}$ at $r = 0$ from
the coincidence weight.  Ansatz $\psi(r) = z^{|r|/2}$ with
$|z| < 1$ gives the bulk eigenvalue $\mu_{\text{rel}} = (1+z)^2/z$
for $r \neq 0$.  The boundary condition at $r = 0$,
$\ee^{\beta^2}\bigl(2z + 2\bigr) = \mu_{\text{rel}}$, yields
$2\ee^{\beta^2}(1+z) = (1+z)^2 / z$, and hence
$z = 1/(2\ee^{\beta^2} - 1)$.  Substituting,
$\mu_{\text{rel}} = 4\ee^{2\beta^2}/(2\ee^{\beta^2}-1)$.
Since $\lambda(1, \beta; \mathbb{Z}) = \beta^2/2 + \log 2$ and
$\lambda(2, \beta; \mathbb{Z}) = \beta^2 + \log 2 + \log \mu_{\text{rel}}
- \log \ee^{\beta^2}$ (the factored-out $\ee^{\beta^2}$ per step and
the COM factor $\log 2$ for $k_s = 0$), a direct computation gives
Eq.~\eqref{eq:d2-closed}.
\end{proof}

\begin{corollary}[Small-$\beta$ asymptotics]
\label{cor:d2-small}
As $\beta \to 0$,
\[
  D(2, \beta; \mathbb{Z}) \;=\; \beta^4 \;-\; \frac{\beta^6}{3}
   \;+\; \frac{\beta^8}{6} \;+\; O(\beta^{10}),
\]
so the leading coefficient is $\beta^4$, matching the
Kardar / Bethe-ansatz scaling of the 2-body bound-state binding for
1D attractive $\delta$-Bose gases.  The bound state has
localisation length $1 / \log(2 \ee^{\beta^2} - 1)$, diverging
as $1/\beta^2$ at weak coupling and shrinking to $O(1)$ lattice
sites at strong coupling.
\end{corollary}

Thm.~\ref{thm:d2-closed} is verified against the numerical
transfer-matrix spectrum: $D(2, \beta; W)$ is monotone decreasing
in $W$ and converges from above to Eq.~\eqref{eq:d2-closed} as
$W \to \infty$.  At $\beta = 0.8$ the closed form gives
$D(2, 0.8; \mathbb{Z}) = 0.25290$, with the finite-$W$ sequence
$0.337, 0.310, 0.294, 0.284, 0.277, 0.272, 0.269$ at
$W = 3, \dots, 9$.  Tests in
\texttt{tests/test\_replica\_closed.py} verify the closed form,
its series expansion coefficient by coefficient, and the finite-$W$
convergence.

\paragraph{Status of the other boundary closed forms.}
The two remaining boundary exponents sit on opposite ends of the
tractability spectrum, and it is worth being explicit about which
is which.

\emph{KPZ upper-tail $c_3$ is closed in continuum.}
The coefficient $c_3 = 4/3$ is known exactly for the continuum
KPZ equation~\citep{DasTsai2021, KrajenbrinkLeDoussal2021,
Tsai2022}, derived in multiple independent ways: each route goes
through the Bethe-ansatz cubic-in-$n$ shape
$\lambda(n) = \tfrac{\beta^2}{2} n + \tfrac{\beta^4}{48} n(n^2-1)$,
followed by a saddle-point / Airy-kernel edge argument on the
resulting Fredholm determinant.  The $3/2$ exponent is the
Legendre conjugate of the cubic.  For \emph{lattice} polymers the
value of $c_3$ is model-dependent; Thm.~\ref{thm:d2-closed} already
captures the $n = 2$ binding on $\mathbb{Z}$ in closed form, and the
analogous three-body coordinate Bethe ansatz on $\mathbb{Z}$
(nested, tractable for on-site attraction) would close the
$n = 3$ sector and give the lattice $c_3$ by the same Legendre
route. This is routine extension of the machinery already in
place rather than a conceptual obstacle.

\emph{LERW $\nu_3$ is genuinely open.}
In $d = 2$ the exponent $\nu_2 = 4/5$ follows from the
conformal-invariance / SLE$_2$ scaling
limit~\citep{LawlerSchrammWerner2004}, with the universal formula
$d_f = 1 + \kappa/8$.  In $d \geq 4$ the exponent $d_f = 2$ comes
from the mean-field regime above the upper critical
dimension~\citep{Lawler1986}.  In $d = 3$ neither mechanism
applies: there is no conformal invariance in the scaling limit and
no small parameter controlling the interaction.  Kozma's 2007
scaling-limit theorem~\citep{Kozma2007} is an \emph{existence}
result --- it shows a limit exists and is supported on a set of
some Hausdorff dimension --- without producing $\nu_3$ in any
form.  No closed form has been proposed even conjecturally.  The
same structural obstruction (below the upper critical dimension,
no conformal symmetry in odd dimension) is what keeps the 3D
Ising critical exponents open.  For our programme this means:
the finite-size Kirchhoff ratio of Prop.~\ref{prop:lerw} is exact
on every finite graph --- the enumeration layer is complete ---
but the scaling limit on $\mathbb{Z}^3$ is beyond what any current
technique can close.

% ==================================================================
\section{Machine verification}
\label{sec:machine}

The three propositions are backed by a passing pytest suite
\texttt{tests/test\_enumerative\_boundary.py}, exercising three new
modules:
\begin{itemize}[leftmargin=*]
\item \texttt{rdft/ac/lerw.py} implements the Kirchhoff ratio of
  Prop.~\ref{prop:lerw} and the length generating function
  Eq.~\eqref{eq:Fab}.  Tests verify $\Tree(K_4) = 16$ (Cayley) and
  $\Tree(C_n) = n$; that LERW probabilities sum to one on $K_4$ for
  several endpoint pairs; that the two arcs of $C_6$ from $0$ to $3$
  carry equal mass $\tfrac12$, while the short and long arcs from
  $0$ to $1$ carry unequal mass summing to one.
\item \texttt{rdft/ac/matrix\_model.py} implements the Harer--Zagier
  recursion for $\varepsilon_g(k)$ and the genus expansion of
  Eq.~\eqref{eq:HZ-GOE} (orientable / GUE case).  Tests verify
  $\varepsilon_0(k) = C_k$ (Catalan), the classical values
  $\varepsilon_1(2,3,4,5) = (1, 10, 70, 420)$, the identity
  $\sum_g \varepsilon_g(k) = (2k-1)!!$ at $N = 1$, and Wigner
  convergence $\Ex\Tr M^{2k} / N^{k+1} \to C_k$ as $N \to \infty$.
\item \texttt{rdft/ac/replica.py} implements direct enumeration of
  $\Ex[Z_N^n]$ from Eq.~\eqref{eq:replica-moment} and a Monte Carlo
  cross-check by disorder averaging.  Tests verify the $n = 1$
  enumeration matches the transfer-matrix path count at $\beta = 0$,
  and that exact enumeration agrees with Monte Carlo within
  $3\sigma$ at $n = 1$ and $n = 2$ for $N = 3$, $\beta = 0.5$.
\end{itemize}

\noindent
These are sanity checks for the combinatorial identities, not
predictions of the open exponents.

% ==================================================================
\section{Expressivity experiments}
\label{sec:experiments}

A second pass of modules pushes each case past the baseline check of
\S\ref{sec:machine} into a non-trivial numerical experiment.

\paragraph{LERW sampling on $\mathbb{Z}^d$ tori
(\texttt{rdft/ac/lerw\_scaling.py}).}
Lawler's loop-erased random-walk construction samples the measure of
Prop.~\ref{prop:lerw} directly: run SRW from $a$ until it hits $b$,
apply chronological loop erasure, record the path length. We verify
on $K_4$ that sampled mean length agrees with the exact Kirchhoff
ratio within one SEM, then sweep $\langle |\gamma|\rangle_{0 \to (r,
0)}$ on $\mathbb{Z}^2_{16}$ and $\mathbb{Z}^3_{8}$.  The mean length
is strictly increasing in $r$ in both dimensions.  On a periodic
torus (no killing boundary) the Kenyon slope $d_f = 5/4$ is not
recovered at these finite sizes; the natural setting for Kenyon's
theorem is a box with Dirichlet boundary, handled next.

\paragraph{Dirichlet-box LERW and recovery of Kenyon / Kozma
exponents (\texttt{rdft/ac/lerw\_dirichlet.py}).}
For the killed-walk LERW on $[0, L{-}1]^d$ --- start at the centre,
run SRW with chronological loop erasure, stop the first time the
walker steps outside the box --- the mean path length scales as
$\langle |\gamma| \rangle \asymp L^{d_f}$ with the lattice LERW
fractal dimension~\citep{Kenyon2000LERW, Kozma2007}.  A Monte-Carlo
sweep at modest compute ($1500$--$3000$ samples per $L$) gives:
\begin{center}
{\small
\begin{tabular}{@{}cccl@{}}
\toprule
$d$ & $L$ range & Fit $d_f$ & Ref.\ $d_f$ \\
\midrule
$2$ & $8$--$48$ & $1.28$ & $5/4$ \\
$3$ & $4$--$12$ & $1.67$ & $\approx 1.624$ \\
$4$ & $3$--$6$  & $1.77$ & $2$ (mean field) \\
\bottomrule
\end{tabular}
}
\\[2pt]
{\small References: $d=2$ \citep{Kenyon2000LERW}; $d=3$
\citep{Kozma2007, Wilson2010}.}
\end{center}
The $d = 2$ slope sits $2.4\%$ above the exact $5/4$ at these box
sizes, the $d = 3$ slope sits $2.8\%$ above Kozma's value, and the
$d = 4$ trend toward the mean-field $2$ is visible but finite-size
dominated.  This is the clean version of the numerical face of
Prop.~\ref{prop:lerw}: the Kirchhoff-ratio measure, sampled on the
graph where Kenyon's theorem applies, recovers both known exponents
within percent-level agreement from the CFAC enumeration layer
alone.

\paragraph{Replica transfer matrix
(\texttt{rdft/ac/replica\_transfer.py}).}
The $n$-walker CRN of Prop.~\ref{prop:kpz-replica} is a transfer
matrix on configuration space $\{-W, \dots, W\}^n$ of size
$(2W+1)^n$.  The Perron eigenvalue gives the replica rate
\begin{equation}
  \lambda(n, \beta; W) \;=\;
   \lim_{N \to \infty} \tfrac{1}{N} \log \Ex[Z_N^n].
\end{equation}
Tests verify the free baseline $\lambda(n, 0; W) = n \log(2
\cos(\pi/(2W+2)))$ (open-chain tight-binding), monotonicity in
$\beta$, superlinear scaling $\lambda(n, \beta; W) / n$ increasing in
$n$ for $\beta > 0$ (replica attraction / binding), and exact
agreement between the transfer-matrix partition function and the
direct CRN enumeration of \S\ref{sec:machine} for small $(n, N)$.

\paragraph{Cubic-in-$n$ Bethe-ansatz signature
(\texttt{rdft/ac/replica\_cubic.py}).}
Kardar's exact Bethe-ansatz solution of the continuum
$n$-body attractive $\delta$-Bose gas~\citep{Kardar1987,
CalabreseLeDoussalRosso2010} gives
\begin{equation}
  \lambda(n, \beta) \;=\; a(\beta)\, n \;+\; b(\beta)\, \frac{n(n^2 - 1)}{6}
  \;+\; O(n^5),
\end{equation}
the Lieb--Liniger ground-state shape that underlies the KPZ
free-energy cumulants.  Define $D(n) := \lambda(n, \beta; W) - n\,
\lambda(1, \beta; W)$.  The pure cubic ansatz predicts
$D(2) = b$ and $D(3) = 4b$, so the model-independent witness
$R(\beta, W) := D(3) / D(2) \to 4$ in the continuum limit
$W \to \infty$.  Exact transfer-matrix evaluation yields:

\medskip
\noindent
\begin{tabular}{@{}rrrrr@{}}
\toprule
$W$ & $D(2)$ & $D(3)$ & $R = D(3)/D(2)$ & $|R - 4|$ \\
\midrule
3 & 0.337 & 1.100 & 3.261 & 0.739 \\
4 & 0.310 & 1.040 & 3.352 & 0.648 \\
5 & 0.294 & 1.006 & 3.420 & 0.580 \\
6 & 0.284 & 0.985 & 3.470 & 0.530 \\
7 & 0.277 & 0.972 & 3.507 & 0.493 \\
8 & 0.272 & 0.962 & 3.535 & 0.465 \\
9 & 0.269 & 0.956 & 3.556 & 0.444 \\
\bottomrule
\end{tabular}

\medskip
\noindent (all at $\beta = 0.8$, exact eigenvalues).  The ratio
approaches $4$ monotonically with the lattice width, with the
finite-$W$ correction empirically consistent with $|R - 4| \sim 1/W$
in the regime tested.  Tests verify monotonicity in $W$,
positivity of binding for $\beta > 0$, vanishing at $\beta = 0$,
and that the lattice ratio sits strictly in $(3, 4)$ at every $W$
in the tested range.  This is the cleanest analytic handle on
$\lambda(n, \beta; W)$ available from the transfer matrix at
small $n$ alone --- the higher-$n$ continuation that would close
the upper-tail coefficient is the next analytic step.

\paragraph{Planar semicircle resolvent and Catalan reconstruction
(\texttt{rdft/ac/tap\_complexity.py}).}
The map-count generating function of Prop.~\ref{prop:tap} is tested
at the planar order: the closed form of
\begin{equation}
  I(f) \;=\; \int_{-2}^{2} \log|\lambda - f|\, \rho_{\rm sc}(\lambda)
  \,\dd\lambda
\end{equation}
is matched against two independent computations.  (i) Direct
numerical integration of the defining integral, at a range of probe
values $f \in \{-3, -2.5, -1.5, 0, 1.2, 2.5, 4\}$; agreement within
$2 \times 10^{-3}$.  (ii) Series reconstruction from Catalan numbers
via
\[
  I(f) \;=\; \log|f| - \sum_{k \geq 1} \frac{C_k}{2k}\, f^{-2k},
  \qquad |f| > 2,
\]
matched within $10^{-4}$ at $|f| \in \{2.5, 3, 4, 6\}$ with $80$
Catalan terms.  This is the genus-$0$ face of Prop.~\ref{prop:tap}:
Catalan counts (planar rooted one-face maps) reconstruct the
semicircle resolvent.

\paragraph{Combinatorial-animal experiment for $\boldsymbol{d=3}$
LERW (\texttt{rdft/ac/lerw\_hierarchical.py}).}
Since the $\mathbb{Z}^3$ thermodynamic generating function is not
$D$-finite, standard analytic-combinatorics extraction does not
terminate.  The alternative route is to build a \emph{combinatorial
animal}: a self-similar graph sequence $\{G_k\}$ on which LERW is
tractable, and measure the gap between the animal's scaling
exponent and the $\mathbb{Z}^3$ Monte-Carlo value~$d_f \approx
1.624$.  Two animals are implemented:

\begin{description}
\item[Migdal--Kadanoff diamond $(b=4, s=2)$.]  $G_{k+1}$ replaces
  each edge of $G_k$ by $b$ parallel branches of $s$ edges each.
  Effective dimension $d_{\text{eff}} = 1 + \log b/\log s = 3$.
  Every $G_k$ is a series-parallel graph, so its UST
  deterministically selects one branch per diamond at every level
  and the source-to-target path realises exactly the graph
  diameter $s^k$.  Hence $d_f^{\text{MK}}(4, 2) = 1$ exactly, a
  large-gap \emph{negative result}: the MK approximant, despite
  having the correct $d_{\text{eff}} = 3$, loses the meandering
  structure of LERW entirely.  The failure is structural: LERW is
  degenerate on every series-parallel graph.

\item[Sierpi\'nski gasket.]  $G_{k+1}$ is the ``triforce''
  arrangement of three copies of $G_k$ glued at top-level corners.
  Hausdorff dimension $\log 3 / \log 2 \approx 1.585$.  The gasket
  is not series-parallel: every sub-triangle carries a cycle, so
  loop erasure is active.  Monte-Carlo extraction at levels
  $k = 2, \dots, 5$ gives $d_f^{\text{Sierp}} \approx 1.20$,
  strictly between $1$ and the gasket's Hausdorff dimension.  The
  animal now has nontrivial loop-erasure exponent but at the
  wrong fractal dimension; its gap to $\mathbb{Z}^3$'s value
  $1.624$ is $\approx 0.42$.
\end{description}

\medskip
\noindent
{\small
\begin{tabular}{@{}p{0.34\columnwidth}ccc@{}}
\toprule
Animal & $d_{\text{eff}}$ & $d_f^{\text{LERW}}$ & Gap \\
\midrule
MK diamond $(4, 2)$  & $3$         & $1.00$         & $-0.62$ \\
Sierpi\'nski gasket  & $\log_2 3$  & $\approx 1.20$ & $-0.42$ \\
$\mathbb{Z}^3_L$ box & $3$         & $\approx 1.67$ & $+0.05$ \\
\bottomrule
\end{tabular}
}

\medskip
\noindent The two animals break in qualitatively different ways.
The MK diamond reproduces the \emph{dimension label} of
$\mathbb{Z}^3$ but not its graph structure (every loop closes
trivially).  The Sierpi\'nski gasket reproduces the
\emph{loop-erasure mechanism} but at the wrong fractal dimension.
$\mathbb{Z}^3$ demands both non-trivial loop structure and
translation-invariant three-dimensional connectivity at once ---
neither animal captures both, which is an honest statement of
what the obstruction to a closed-form $\mathbb{Z}^3$ exponent
actually is.

\paragraph{Algebraic / compositional reading.}
Each hierarchical construction $G_{k+1} = \Phi(G_k)$ is a finite
\emph{compositional} rewrite in the sense of Flajolet--Sedgewick's
symbolic method: at the level of generating functions, the LERW
length series $F_k(z)$ satisfies a functional recursion
\begin{equation}
  F_{k+1}(z) \;=\; \Phi^\sharp \bigl[ F_k(z) \bigr],
  \label{eq:rg-recursion}
\end{equation}
where $\Phi^\sharp$ is the rational transformation on
generating functions induced by the graph rewrite $\Phi$ through
the Kirchhoff-ratio definition of Prop.~\ref{prop:lerw}.  The
fractal exponent $d_f$ is then the linearisation eigenvalue of
$\Phi^\sharp$ at its renormalisation fixed point:
\[
  d_f \;=\; \frac{\log \lambda^\sharp}{\log s},
\]
with $\lambda^\sharp$ the dominant eigenvalue of
$\mathrm{d}\Phi^\sharp$ and $s$ the length-scale factor of
$\Phi$.  This is the Kadanoff block-spin renormalisation rephrased
as a fixed-point problem in the Flajolet symbolic category: the
\emph{algebra of constructions} (SEQ, PROD, recursive
substitution) determines the \emph{algebra of generating
functions}, and the exponent comes from a single eigenvalue
derivative.  For series-parallel $\Phi$ the transformation is
degenerate and $\lambda^\sharp = s$ (hence $d_f = 1$); for
non-series-parallel $\Phi$ the eigenvalue carries non-trivial
graph information.  The obstruction at $\mathbb{Z}^d$ is precisely
that there is no finite rewrite $\Phi$ that generates $\mathbb{Z}^d$
by self-similar substitution: translation invariance is a
\emph{limit} construction rather than a compositional one, so the
symbolic-method machinery has no finite animal to seed.

\paragraph{Two further animal experiments and their negative
findings.}
Having exhausted the self-similar hierarchical approach, the
natural next question is whether a tower of
\emph{non-self-similar} but still algebraic animals can close in
on $d_f^{(3)}$.  Two are implemented and documented here for
completeness; both fail, and the failure mode of each is
informative.

\emph{(a) Finite-cross-section tubes
(\texttt{rdft/ac/lerw\_tube.py}).}
LERW on the tube $\{0, \dots, N-1\}^2_{\text{periodic}} \times
\{0, \dots, L-1\}_{\text{Dirichlet}}$ defines, for each $N$, an
$N^2$-state Markov chain with an algebraic spectrum, hence an
algebraic candidate exponent $d_f^{(N)}$.  Monte-Carlo extraction
at $N = 1, 2, 3, 4, 5$ with $L \in \{8, 12, 16, 24, 32\}$ gives
$d_f^{(N)} \approx 1$ for every $N$ tested (ballistic collapse).
The reason is structural: a narrow tube is quasi-1D, and LERW on
a 1D path is a deterministic monotone traversal.  Cross-sectional
wandering contributes a sub-leading constant factor, not a
different power of $L$.  The tower of tube exponents does not
approach $d_f^{(3)}$ as $N$ grows --- it stays at $1$ until
$N \sim L$, at which point the ``tube'' is no longer a tube and
one is simply looking at an isotropic box.  This rules out the
tube-transfer-matrix approach as a tower of algebraic animals
for $d_f^{(3)}$.

\emph{(b) Algebraic extrapolation on isotropic boxes
(\texttt{rdft/ac/lerw\_extrap.py}).}
The correct target observable is the isotropic-box mean length
$\langle|\gamma|\rangle_L$ on $\{0, \dots, L-1\}^3$, which our
Dirichlet sampler produces.  Classical AC-style algebraic
extrapolation applied to the sequence
$\{\langle|\gamma|\rangle_L\}$ includes:
\begin{itemize}[leftmargin=*]
\item Naive log-log fit;
\item $(1 + B/L^\omega)$ correction-to-scaling fits with $\omega$
  fixed or grid-searched;
\item Richardson extrapolation on effective exponents
  $d_f^{\text{eff}}(L) := \log(\langle|\gamma|\rangle_{L+}/
  \langle|\gamma|\rangle_{L-}) / \log(L_+ / L_-)$, assuming a
  $1/L$ leading correction;
\item Neville--Richardson tableau extrapolation, assuming an
  integer-power correction tower $1/L, 1/L^2, \dots$
\end{itemize}
On clean synthetic data with a known correction structure the
schemes recover the injected exponent (tests witness this).  On
real LERW data in $d = 2$ and $d = 3$, however, the
``sophisticated'' schemes do \emph{not} outperform the naive fit:
they systematically bias $d_f$ downward because the true
correction to scaling is not a simple power of $1/L$.  In $d = 3$
at $L \in \{4, 6, \dots, 16\}$ the naive fit gives $d_f \approx
1.62$ (coincidentally within a few parts per thousand of Wilson's
$1.6236$), while the $(1 + B/L)$ correction fit drags it to
$\approx 1.54$ and Richardson extrapolation to $\approx 1.29$.
This is a concrete epistemic statement: \emph{LERW finite-size
corrections do not live in the algebraic-power-of-$1/L$ category
that Richardson and Neville assume.}  If they did, those schemes
would dominate the naive fit by orders of magnitude.  They don't,
so the corrections include some non-algebraic structure
(logarithmic factors, non-integer powers, or boundary-specific
terms that do not obey the $\Phi^\sharp$-eigenvalue story).

\paragraph{A fifth animal: the sandpile group
(\texttt{rdft/ac/sandpile\_group.py}).}
The Majumdar--Dhar bijection~\citep{MajumdarDhar1992} maps
recurrent configurations of the abelian sandpile model (ASM) on a
graph $G$ to spanning trees of $G$; Wilson's theorem then maps
the $a$-to-sink path in the UST to LERW.  So ASM states, spanning
trees, and LERW carry the same probabilistic content.  The
algebraic upgrade is that recurrent ASM configurations form a
finite abelian group --- the \emph{sandpile group}
\[
  K(G) \;=\; \mathbb{Z}^{V \setminus \{s\}}
   \,/\, \Delta'\, \mathbb{Z}^{V \setminus \{s\}},
\]
with $\Delta'$ the Laplacian with the sink row/column removed.
The sequence $\{K(\mathbb{Z}^d_L)\}_{L=1,2,\dots}$ is a concrete
tower of finite algebraic animals, each of which is completely
determined by the Smith normal form of $\Delta'$.

Computation (via integer Smith normal form): for
$\mathbb{Z}^2_L$ the non-trivial invariant factors are
$L = 3$: $[6, 6, 18, 18]$;
$L = 4$: $[2, 2, 8, 24, 24, 24, 96]$;
$L = 5$: $[10, 10, 50, 50, 50, 50, 50, 50]$;
for $\mathbb{Z}^3_3$ torus,
$[6, 6, 6, 18, 18, 54, 54, 54, 54, 162, 162, 162]$.
The group orders (= Matrix-Tree spanning-tree counts) are verified
against the closed-form Laplacian-eigenvalue product, so the
computation is correct.

\emph{Fifth failure mode.}
The group structure carries clear number-theoretic content (the
divisors are products of prime powers determined by $L$'s
factorisation), but its scaling is erratic rather than universal:
the largest invariant factor grows as $18, 96, 50, \dots$ for
$L = 3, 4, 5$, giving effective exponents
$\log(d_{\max}) / \log L \in \{2.63, 3.29, 2.43\}$ --- dominated
by number-theoretic fluctuations rather than any clean power-law.
The obstruction is structural: the sandpile GROUP is the Smith
normal form of $\Delta'$, a \emph{static} algebraic invariant.
The LERW \emph{fractal dimension} is a dynamical observable --- the
expected length of the UST path between two points --- that
depends on how the full Laplacian spectrum interacts with the
path decomposition, not on the Smith-normal-form decomposition
alone.  What WOULD carry $d_f$ is the avalanche-size / wave-profile
distribution of the ASM \emph{dynamics}; that is not an algebraic
invariant of $K(G)$.

The upshot: the Majumdar--Dhar bijection is real and deep, and
the sandpile group is a legitimate algebraic animal for every
finite graph, but it is the wrong algebraic object to extract
$d_f^{(3)}$ from.  The group's content is Kirchhoff (spanning-tree
count, hence bulk entropy); it does not carry path-length
scaling.  This is another honest no-go to add to the catalogue.

\paragraph{What IS the right combinatorial object for path
length? (\texttt{rdft/ac/lerw\_exact.py})}
With five algebraic animals catalogued as failures, a sharper
question is: \emph{what is the combinatorial object that actually
captures dynamical path length, the thing we ought to have been
extracting asymptotics from?}  The answer has been in the paper
from the start: it is the length generating function of Prop 1
itself,
\begin{equation}
  F_{ab}(z; G) \;=\; \sum_{\gamma : a \to b\,\text{simple}}
    z^{|\gamma|}\; \frac{w(\gamma)\, \Tree(G / \gamma)}{\Tree(G)}.
  \label{eq:F-path}
\end{equation}
$F_{ab}(z; G)$ is a \emph{rational function in $z$} on every finite
graph $G$, with mean length $\langle|\gamma|\rangle = F'_{ab}(1)$.
This is deterministic --- no Monte Carlo --- and is the exact
algebraic object whose asymptotic behaviour in $L$ would give
$d_f^{(d)}$.

Two sanity checks.  Naive hope: maybe $\langle|\gamma|\rangle$
equals the sum of absolute electric currents
$\sum_e |i_e(a \to b)|$, where $i_e$ is the unit current on edge
$e$ in the resistor network with unit edge conductances.  This is
true on trees.  On $K_4$ from vertex $0$ to $1$ the sum of
absolute currents is $\tfrac12 + 4 \cdot \tfrac14 + 0 = \tfrac32$
while the exact mean from Eq.~\eqref{eq:F-path} is $13/8 = 1.625$
--- the discrepancy comes from edge $(2, 3)$, which carries zero
current by symmetry but still appears in $1/8$ of LERW paths.
So the correct combinatorial object is \emph{not} a sum of currents;
it is the Kirchhoff ratio summed against $z^{|\gamma|}$ over simple
paths, exactly as in Eq.~\eqref{eq:F-path}.

Deterministic values from Eq.~\eqref{eq:F-path} via direct path
enumeration plus Kirchhoff determinants:
\begin{center}
{\small
\begin{tabular}{@{}lcc@{}}
\toprule
Case (corner-to-corner) & Exact $\langle|\gamma|\rangle$ & MC check \\
\midrule
$\mathbb{Z}^2_3$ & $17/4 = 4.25$          & $4.26 \pm 0.01$ \\
$\mathbb{Z}^2_4$ & $6.69228\ldots$         & $6.69 \pm 0.02$ \\
$\mathbb{Z}^2_5$ & $9.29526\ldots$         & $9.30 \pm 0.02$ \\
$\mathbb{Z}^3_2$ & $53/16 = 3.3125$        & $3.31 \pm 0.02$ \\
$\mathbb{Z}^3_3$ & $7.84146\ldots$ ($\sim 2$ min enum.)
                                           & $7.85 \pm 0.05$ \\
\bottomrule
\end{tabular}
}
\end{center}

These are \emph{exact rationals} (the last two decimalised for
brevity), and they match Monte Carlo to sampling precision.
The length generating function is therefore correctly identified
as the combinatorial object for path length.

\emph{Why we didn't already extract $d_f^{(3)}$ from it.}
Two obstructions --- the same two that shaped the whole paper:
\begin{itemize}[leftmargin=*]
\item \textbf{Computation cost.}  The number of simple paths on
  $\mathbb{Z}^3_L$ grows super-exponentially in $L$; direct
  enumeration becomes infeasible beyond $L = 3$.  No known
  analytic shortcut reduces the Kirchhoff-ratio sum to anything
  D-finite as $L \to \infty$ --- because the scaling limit IS the
  non-D-finite lattice Green's function.
\item \textbf{Structural obstruction.}  Substituting
  $z = 1$ in Eq.~\eqref{eq:F-path} gives $1$ (normalisation); the
  derivative at $1$ is a \emph{weighted resistance-like quantity}
  against the full Laplacian spectrum, and extracting its
  $L$-scaling is exactly the closed-form problem we cannot solve
  in $d = 3$.
\end{itemize}

So the right combinatorial object is in our hand; the
rate-limiting step is not identifying it but extracting
its asymptotics at scale.  The five algebraic animals catalogued
above are the known attempts at shortcut asymptotics, each
losing a different ingredient of $F_{ab}(z; \mathbb{Z}^d_L)$ in
the process.  A wider AC category (resurgent transseries,
non-D-finite closure, determinantal / Fredholm approximation of
the scaling limit) is what would be required to close the
extraction --- and that is the honest open direction this paper
hands off.

\paragraph{Notebook: a special deep problem.}
The two negative results above, taken with the MK diamond,
Sierpi\'nski gasket, and sandpile-group animals, catalogue five
attempted algebraic animals for $d_f^{(3)}$:
\begin{center}
{\small
\begin{tabular}{@{}p{0.30\columnwidth}p{0.58\columnwidth}@{}}
\toprule
Animal & Failure mode \\
\midrule
MK diamond $(4, 2)$
 & Correct $d_\text{eff} = 3$ but series-parallel; UST is
   trivially ballistic. $d_f = 1$. \\
Sierpi\'nski gasket
 & Correct loop-erasure structure but wrong fractal dimension
   ($d_H = \log_2 3$, not 3). $d_f \approx 1.20$. \\
$\mathbb{Z}^3$ tube tower
 & Quasi-1D at every finite $N$. $d_f^{(N)} = 1$ for all
   accessible $N$. \\
Algebraic extrapolation on Dirichlet boxes
 & Naive fit hits $\sim 1.62$, but no scheme improves on it ---
   corrections are not a clean algebraic power of $1/L$. \\
Sandpile-group tower (Majumdar--Dhar bijection)
 & Right bijection to LERW; group carries only Matrix-Tree
   content (spanning-tree count), not path-length scaling.
   Erratic $\log d_{\max}/\log L$. \\
\bottomrule
\end{tabular}
}
\end{center}

\noindent The five failures have a common theme: they each succeed
at capturing \emph{one} of the ingredients of $\mathbb{Z}^3$
(dimension, loops, finite Markov structure, algebraic corrections,
or the Majumdar--Dhar algebraic-group upgrade) while losing the
dynamical path-length content that $d_f$ actually measures.  The
target object sits at the intersection of all five, and no finite
algebraic construction we have been able to put to it captures the
intersection.

This is a concrete statement of the ``deep'' part of the problem
--- stronger than ``we don't have a closed form''.  It says: the
generating-function category within which we can do AC asymptotic
extraction is populated by objects, none of which contains the
$\mathbb{Z}^3$ LERW two-point function in its closure under
composition.  Whether a \emph{wider} AC category (resurgent
transseries, hyperasymptotic Borel summation, non-D-finite
closure) contains it is the honest open question this paper
leaves open.

\paragraph{What the experiments do and do not demonstrate.}
Each module now has a concrete computational handle tested against
an independent reference:
\begin{itemize}[leftmargin=*]
\item LERW: samplers consistent with the exact Kirchhoff ratio, and
  numerical recovery of Kenyon's $d_f = 5/4$ in $d=2$ and Kozma's
  $d_f \approx 1.624$ in $d=3$ from a Dirichlet-box Lawler sweep
  (each within $\sim 3\%$ at modest sample size); plus the
  combinatorial-animal experiment above, which quantifies exactly
  what self-similar surrogates capture and miss about $\mathbb{Z}^3$.
\item KPZ: a transfer-matrix replica partition function consistent
  with the CRN enumeration and with the free-walker spectrum; the
  exact $\lambda(n, \beta; W)$ curve at $n = 1, 2, 3$ shows the
  Bethe-ansatz cubic-in-$n$ shape, with $D(3)/D(2) \to 4$
  (Kardar continuum limit) monotonically as $W$ grows; and the
  closed-form 2-body replica rate
  $D(2, \beta; \mathbb{Z}) = 2\beta^2 - \log(2\ee^{\beta^2} - 1)$
  (Thm.~\ref{thm:d2-closed}) is the Tier-A deliverable for the
  $n = 2$ sector.
\item TAP: the Catalan-series reconstruction of the semicircle
  resolvent works to the precision of the truncation.
\end{itemize}
The open exponents and coefficients are still open.  What these
experiments show is that the CFAC enumeration layer produces
concrete, testable objects in every case --- not just structural
identities.

% ==================================================================
\section{What the enumeration layer buys}
\label{sec:buys}

Prop.~\ref{prop:lerw}, \ref{prop:tap}, and~\ref{prop:kpz-replica}
are not new mathematics.  The point is that they share a common
CFAC-enumerative structure:

\medskip
\noindent
\begin{tabular}{@{}p{0.28\columnwidth}@{\hspace{4pt}}p{0.64\columnwidth}@{}}
\toprule
Case & Enumerated object \\
\midrule
LERW 3D       & Simple paths weighted by $\Tree(G/\gamma)/\Tree(G)$ \\
TAP complexity & Rooted one-face maps by genus (Harer--Zagier) \\
KPZ moments   & $n$-tuples of walkers by coincidence count \\
\bottomrule
\end{tabular}
\medskip

Each is a weighted sum over a discrete combinatorial set, with
weights that are either first Symanzik polynomials
(Eq.~\eqref{eq:kirchhoff}), ribbon-graph Euler numbers
(Eq.~\eqref{eq:HZ-GOE}), or Doi--Peliti vertex factors
(Eq.~\eqref{eq:nwalker}).  The existing CFAC pipeline handles all
three.

\paragraph{What the enumeration layer does \emph{not} buy.}
\begin{itemize}[leftmargin=*]
\item For LERW, the 3D exponent $\nu_3$ is still a hard scaling
  limit of the finite-graph formula; no closed form falls out.
\item For TAP, the genus expansion organises finite-$N$ corrections
  but does not by itself give the complexity density $\Phi_0$ in
  closed form for general mixed $p$-spin models.
\item For KPZ, integer-$n$ replica enumeration is a classical
  weighted walker count; the open problem (upper-tail coefficient
  universality) lives at the $n \to 0$ continuation, outside the
  enumeration.
\end{itemize}

What it does buy is a concrete combinatorial definition of the object
under study, with finite-size exactness, and a clean statement of
where the stratification theorem would have to be extended (scalar
$\to$ matrix bridge; algebraic $\to$ non-$D$-finite GF;
deterministic $\to$ disorder-averaged vertex).

\paragraph{Where this leaves the scope map.}
Combining this note with~\citep[\S7, \S8]{AmarteifioMannaC3} gives a
two-tier scope statement:

\begin{slotbox}
\emph{Tier A (stratification theorem)}: polynomial DSEs with
algebraic / $D$-finite generating function and a scalar bridge.
Delivers a skeleton exponent $\tau_0 = 1 + 1/k$ plus a computable
one-loop dressing $\gamma_k(d)$.

\emph{Tier B (enumeration)}: any weighted combinatorial object
(Feynman diagrams, spanning trees, ribbon graphs, $n$-walker paths,
partition sums) with Symanzik / Laplacian / Schur structure.
Delivers exact finite-size generating functions and explicit rational
/ determinantal representations.
\end{slotbox}

The three boundary cases of~\citep{AmarteifioMannaC3} fall outside
Tier A but strictly inside Tier B.  The scope-claim
``CRN-in-field is the underlying object'' in~\citep[\S8]{AmarteifioMannaC3}
is honest at Tier B; the stronger claim that the stratification
theorem applies is honest only within Tier A.

% ==================================================================
\section{Discussion}
\label{sec:discussion}

The three propositions above are classical results rewritten in
CFAC enumeration language.  The value of the rewrite is that it
makes the obstruction for each case \emph{localisable}: LERW 3D
needs non-$D$-finite resummation of a Kirchhoff ratio; TAP complexity
needs a matrix-valued extension of the scalar bridge to close the
complexity density into closed form for general $p$-spin;
KPZ upper tails need a disorder-averaged bridge to fix the
coefficient from the Schur edge.

Two further cases that the companion paper lists only in passing ---
BRST-quantised gauge theory and replica $n \to 0$ --- fit the same
template.  BRST cohomology is a graded enumeration of ghost-number
sectors; replica $n \to 0$ is a continuation of an enumeration.
Neither is a Tier-A object.  Both are, at the enumeration level,
concrete combinatorial constructions the pipeline can handle at
fixed $n$ or fixed ghost number.

\paragraph{Next steps.}
The three modules of \S\ref{sec:machine} establish the baseline;
the interesting extensions are:
(i) push the LERW Kirchhoff-ratio computation to
$\mathbb{Z}^3_N$ for $N \lesssim 20$ via a Fourier-diagonalised
Laplacian, and track the finite-size scaling of
$\partial_z \log F_{0b}$ against the Kozma exponent;
(ii) extend the genus-graded ribbon-graph enumeration to GOE
(non-orientable surfaces, Goulden--Jackson coefficients) and
plug it into the Kac--Rice Legendre transform to recover the
$p$-spin complexity density as a map-count resummation;
(iii) implement the Schur-process reduction Eq.~\eqref{eq:schur-form}
for integrable polymer geometries and verify it against the
direct CRN enumeration on small instances
(e.g.\ $N \leq 4$, $n \leq 3$).

None of these closes the open exponent or coefficient.  Each
sharpens the line between what CFAC-as-theorem covers and what
CFAC-as-enumeration covers, and supplies a concrete object for
downstream work.

% ==================================================================
\begin{thebibliography}{99}
\small

\bibitem[Amarteifio(2026a)]{AmarteifioCFAC2026}
  Amarteifio, S., 2026.
  \textit{CFAC theorem: stratification of Dyson--Schwinger equations
  by Puiseux order.}  Percolation Labs preprint.

\bibitem[Amarteifio(2026b)]{AmarteifioMannaC3}
  Amarteifio, S., 2026.
  The $\Cstr{3}$ Stratum in Non-DP Universality: Slotting Manna and
  Conserved-DP, and what currently resists slotting.
  Percolation Labs preprint
  (\texttt{paper/cfac/manna\_c3\_slotting.tex}).

\bibitem[Amarteifio(2026c)]{Amarteifio2026LDWC}
  Amarteifio, S., 2026.
  Reading the LDWC two-loop roughness coefficient through the
  CFAC lens.
  Percolation Labs preprint
  (\texttt{paper/cfac/ldw\_kirchhoff\_recovery.tex}).

\bibitem[Auffinger, Ben Arous and \v{C}ern\'y(2013)]{ABAC2013}
  Auffinger, A., Ben Arous, G.\ and \v{C}ern\'y, J., 2013.
  Random matrices and complexity of spin glasses.
  \textit{Comm. Pure Appl. Math.} \textbf{66}, 165.

\bibitem[Borodin and Corwin(2014)]{BorodinCorwin2014}
  Borodin, A.\ and Corwin, I., 2014.
  Macdonald processes.
  \textit{Probab. Theory Relat. Fields} \textbf{158}, 225.

\bibitem[Corwin and Ghosal(2020)]{CorwinGhosal2020}
  Corwin, I.\ and Ghosal, P., 2020.
  Lower tail of the KPZ equation.
  \textit{Duke Math. J.} \textbf{169}, 1329.

\bibitem[Goulden and Jackson(1997)]{GouldenJackson1997}
  Goulden, I.~P.\ and Jackson, D.~M., 1997.
  Maps in locally orientable surfaces, the double coset algebra,
  and zonal polynomials.
  \textit{Canad. J. Math.} \textbf{49}, 865.

\bibitem[Harer and Zagier(1986)]{HarerZagier1986}
  Harer, J.\ and Zagier, D., 1986.
  The Euler characteristic of the moduli space of curves.
  \textit{Invent. Math.} \textbf{85}, 457.

\bibitem[Das and Tsai(2021)]{DasTsai2021}
  Das, S.\ and Tsai, L.-C., 2021.
  Fractional moments of the stochastic heat equation.
  \textit{Ann. Inst. H. Poincar\'e} \textbf{57}, 778.

\bibitem[Krajenbrink and Le Doussal(2021)]{KrajenbrinkLeDoussal2021}
  Krajenbrink, A.\ and Le Doussal, P., 2021.
  Inverse scattering solution of the weak noise theory of the KPZ
  equation with flat and Brownian initial conditions.
  \textit{Phys. Rev. E} \textbf{105}, 054142.

\bibitem[Tsai(2022)]{Tsai2022}
  Tsai, L.-C., 2022.
  Exact lower-tail large deviations of the KPZ equation.
  \textit{Duke Math. J.} \textbf{171}, 1879.

\bibitem[Lawler(1986)]{Lawler1986}
  Lawler, G.~F., 1986.
  Intersections of random walks in four dimensions. II.
  \textit{Comm. Math. Phys.} \textbf{97}, 583.

\bibitem[Lawler, Schramm and Werner(2004)]{LawlerSchrammWerner2004}
  Lawler, G.~F., Schramm, O.\ and Werner, W., 2004.
  Conformal invariance of planar loop-erased random walks and
  uniform spanning trees.
  \textit{Ann. Probab.} \textbf{32}, 939.

\bibitem[Calabrese, Le Doussal and Rosso(2010)]{CalabreseLeDoussalRosso2010}
  Calabrese, P., Le Doussal, P.\ and Rosso, A., 2010.
  Free-energy distribution of the directed polymer at high
  temperature.
  \textit{EPL} \textbf{90}, 20002.

\bibitem[Majumdar and Dhar(1992)]{MajumdarDhar1992}
  Majumdar, S.~N.\ and Dhar, D., 1992.
  Equivalence between the Abelian sandpile model and the $q \to 0$
  limit of the Potts model.
  \textit{Physica A} \textbf{185}, 129. (Establishes the
  bijection between recurrent sandpile configurations and
  spanning trees.)

\bibitem[Kardar(1987)]{Kardar1987}
  Kardar, M., 1987.
  Replica Bethe ansatz studies of two-dimensional interfaces with
  quenched random impurities.
  \textit{Nucl. Phys. B} \textbf{290}, 582.

\bibitem[Kenyon(2000)]{Kenyon2000LERW}
  Kenyon, R., 2000.
  The asymptotic determinant of the discrete Laplacian.
  \textit{Acta Math.} \textbf{185}, 239. (Gives
  $d_f = 5/4$ for 2D loop-erased random walk.)

\bibitem[Kozma(2007)]{Kozma2007}
  Kozma, G., 2007.
  The scaling limit of loop-erased random walk in three dimensions.
  \textit{Acta Math.} \textbf{199}, 29.

\bibitem[Krajenbrink(2020)]{Krajenbrink}
  Krajenbrink, A., 2020.
  From Painlev\'e transcendents to the KPZ equation.
  Ph.D.\ thesis, Ecole Normale Sup\'erieure, Paris.

\bibitem[Subag(2017)]{Subag2017}
  Subag, E., 2017.
  The complexity of spherical $p$-spin models.
  \textit{Ann. Probab.} \textbf{45}, 3385.

\bibitem[Wilson(1996)]{Wilson1996}
  Wilson, D.~B., 1996.
  Generating random spanning trees more quickly than the cover time.
  \textit{Proc.\ 28th STOC}, 296.

\bibitem[Wilson(2010)]{Wilson2010}
  Wilson, D.~B., 2010.
  The dimension of loop-erased random walk in 3D.
  \textit{Phys. Rev. E} \textbf{82}, 062102.

\end{thebibliography}

\end{document}
